% 《阿拉伯世界的1500年》期末报告
% 编译：双击 compile.bat（仅 xelatex ×2，IEEE 数字引用）
\documentclass[UTF8,a4paper,12pt]{ctexart}

\usepackage{geometry}
\usepackage{setspace}
\usepackage{graphicx}
\usepackage{caption}
\usepackage{hyperref}

\geometry{top=2.54cm,bottom=2.54cm,left=3cm,right=2.2cm}
\setstretch{1.5}
\setlength{\parindent}{2em}
\captionsetup{font=small,labelfont=bf}
\hypersetup{colorlinks=true,linkcolor=black,citecolor=blue,urlcolor=blue}
\graphicspath{{images/}}
% 限制图片高度，避免大图导致 xelatex 长时间卡住
\newcommand{\reportfig}[2]{%
  \includegraphics[width=#1,height=0.36\textheight,keepaspectratio]{#2}}

\title{\textbf{西西里岛与阿拉伯文化的关系}\\[0.4em]
  \large 《阿拉伯世界的1500年》课程期末报告}
\author{小组合作}
\date{2026年5月27日}

\begin{document}
\maketitle
\thispagestyle{empty}
\vspace{1em}
\begin{center}
  \small 说明：正文汉字约4800字；封面、图注、参考文献与引用标注不计入字数。引用格式为 IEEE 顺序编码制。
\end{center}
\newpage

\section{引言}

中世纪地中海并非文明间的封闭边界，而是持续交换的``液态空间''：钱币、建筑、知识、食物与人口在此流动，形成超越单一政体与单一宗教的混合秩序。课程《阿拉伯世界的1500年》强调，帝国长治久安往往不只靠刀剑，更依赖文书、称号与象征资源的制度化运用；诺曼西西里正是这一命题在地中海西部的典型样本。本文选取``西西里岛与阿拉伯文化的关系''为题，以威尼斯杜卡特金币为物质史切入点，结合诺曼时期的行政遗产、建筑艺术、日常生活与翻译实践，讨论9--13世纪阿拉伯世界与拉丁基督教世界的深层互动。

从学术脉络看，传统叙事常把西西里简化为``拜占庭—阿拉伯—诺曼''的三段更替，却较少追问：政权更替之后，技术、符号与生活方式为何仍能延续？与``文明冲突''范式不同，本文主张跨文化互动常呈现为\textbf{底层制度延续}与\textbf{顶层符号重组}的分层策略\cite{johns2002,metcalfe2009,grierson1990}。诺曼统治者继承迪旺（\emph{dīwān}）档案、土地丈量与税制，却在铸币、宫廷礼仪与对外联姻中嵌入拉丁基督教符号；威尼斯杜卡特则以圣像与拉丁铭文将神学—政治话语铸入流通\cite{stahl2000,travaini2003}。二者并置说明：混合文明是地中海多元社会面对治理难题时的制度性回应，而非征服之后的偶然副产品。下文依钱币—地理—建筑—日常—知识五条线索展开，以期在课程要求的4000--5500字篇幅内做到史料与论点并重。

\section{威尼斯杜卡特：视觉政治经济学}

\begin{figure}[htbp]
  \centering
  \reportfig{0.42\textwidth}{image1.jpeg}
  \caption{威尼斯泽奇诺—杜卡特（\emph{zecchino/ducat}）背面：基督立于曼多拉（\emph{mandorla}）内，外环拉丁铭文将公国奉献于基督，体现``钱币即宣传''（\emph{coinage as propaganda}）的传统。}
  \label{fig:ducat}
\end{figure}

选择杜卡特作为开篇，还因它与课程第12讲所讨论的奥斯曼合法性工具有结构相似性：二者皆把政治正当性写入可复制的媒介——奥斯曼以称号、建筑与花押（\emph{tughra}）宣示统治，威尼斯与诺曼西西里则分别以金币图像与宫廷织物铭文达成类似功能\cite{hitip1958,stahl2000}。图\ref{fig:ducat}所示金币属13世纪以降趋于稳定的威尼斯杜卡特（\emph{zecchino/ducat}）传统，在东地中海贸易网络中获得广泛信用认同\cite{stahl2000,grierson1990}。币面中央为基督像，头戴十字光轮，右手降福、左手持福音书；外环拉丁铭文作 \emph{Sit tibi, Christe, datus, quem tu regis, iste ducatus}（``基督啊，将此公国献于你所治者''），将统治正当性压缩于支付媒介\cite{travaini2003}。杜卡特之成为``硬通货''，并不只因含金量稳定，更因图像—铭文程序在跨区域流通中被不断识别与重复确认：商人、税吏与教会人士在交接金币时，同时完成了对威尼斯共和国神授治理叙事的再确认。

与之对照，伊斯兰铸币长期以阿拉伯书法、宗教名号（如 \emph{lā ilāha illā Allāh}）及几何、植物纹样为主，在公共媒介上回避以人像承担神圣或政治崇拜功能\cite{grierson1960,heidemann2010,hillenbrand1999}。这一传统在倭马亚、阿拔斯以来的金币范式中被制度化，并与伊斯兰教对偶像（\emph{ṣanam}）的伦理讨论相互强化。拉丁世界（尤其受拜占庭图像学影响的意大利城邦）则习惯在最高面值贵金属上安置可见圣像与环绕铭文，使``何者被看见、何者被诵读''同步进入经济交换与政治合法性建构\cite{grierson1990,stahl2000}。杜卡特因而可视为\textbf{视觉政治经济学}：它在港口、市集与账簿之间流动，却把神学—政治话语铸入每一次称重、交割与信任评估的感官经验之中。

对诺曼西西里研究而言，威尼斯杜卡特宜作\textbf{对照类型}，而非本岛铸币的直接标本。考古与文献表明，伊斯兰治下及过渡时期在西西里流通的 \emph{tarì}（tarì）仍以书法与名号为主；诺曼君主国在整合地中海贸易时，必须同时与威尼斯、拜占庭及西方基督教世界的图像—铭文货币文化对话\cite{johns2002,metcalfe2009,metcalf1998,loud2012}。罗杰二世及其后继者并未彻底废除阿拉伯式货币记忆，而是在国家礼仪、加冕服饰与建筑象征中重组``谁有权书写合法性''的问题。核心追问因此不是``欧洲钱币是否更精美''，而是：当迪旺文书、土地册封与税制被继承时，最高层公共媒介的图像程序是否同样被继承，抑或被选择性替换？钱币史在此与课程所讲的奥斯曼``称号、建筑、花押''合法性语言形成跨时空呼应：帝国秩序往往通过可重复的视觉—文本装置得以维系，而非仅依赖军事胜利本身。

\section{地理枢纽与王朝更迭}

\begin{figure}[htbp]
  \centering
  \reportfig{0.72\textwidth}{image2.jpeg}
  \caption{巴勒莫诺曼宫（\emph{Palazzo dei Normanni}）外观：诺曼式体量与阿拉伯式连拱、盲拱装饰并存，是``诺曼—阿拉伯—拜占庭''建筑合成的典型地标。}
  \label{fig:palace}
\end{figure}

西西里面积约2.55万平方公里，是地中海最大岛屿之一：东北经墨西拿海峡与亚平宁半岛相隔仅约3公里，西南与北非突尼斯隔海相望\cite{sicilygeo2026}。在缺乏现代运河的古代，绕行西西里几乎是东西向航线的必经节点；控制该岛，意味着控制贸易税、补给港与舰队机动空间\cite{finley1986,alshaar2026}。岛内地形以山地丘陵为主，埃特纳火山喷发形成的火山灰土壤极为肥沃，利于小麦、橄榄与葡萄种植；沿海良港则便于舰队停泊。公元827年起，阿拔斯王朝支持的远征军逐步征服西西里，831年占领巴勒莫后，将其由拜占庭边陲前哨改造为伊斯兰西西里的政治与经济中心。

从更长时段观察，西西里先后经历希腊殖民、罗马行省、拜占庭统治与阿拉伯埃米尔国（827--1091年）及诺曼王国（1130年后）等阶段；每一次政权更替都留下层次累加的文化沉积。阿拉伯埃米尔国定都巴勒莫后，兴建清真寺、宫殿与集市，并完善港口管理，使该城首次具备与巴格达、开罗等大都会相媲美的城市规模与国际联系。这对理解阿拉伯文化向西地中海扩散具有关节意义。阿拉伯文化进入西西里，是地理诱惑、帝国扩张与拜占庭内乱（如西西里总督叛乱）交织的结果；但文化影响力的曲线并不与军事统治完全重合。11世纪末诺曼人完成征服后，穆斯林政治权力退场，伊斯兰文化元素却在诺曼宫廷庇护下获得更系统的整合\cite{metcalfe2009,finley1986}——巴勒莫由此成为``诺曼—阿拉伯—拜占庭''合成文化的枢纽。Nuha Alshaar 等学者近年研究亦强调，西西里穆斯林遗产在中世纪晚期及近代早期仍以地名、农业技术与社区记忆的形式持续存在\cite{alshaar2026}。这一``征服结束而文化延续''的现象，为理解阿拉伯世界与欧洲的关系提供了关键反例：地中海东部与南部的知识、作物与制度，往往通过海路向西渗透，并在新君主国的框架内被重新命名与再利用。

\section{建筑、艺术与王权符号}

诺曼征服后，阿拉伯装饰语言与空间逻辑并未被简单抹去，而在新王权体系中被策略性吸收，形成学界所称``诺曼—阿拉伯—拜占庭''合成文化\cite{finley1986,hillenbrand1999}。罗杰二世（1130--1154年在位）尤其以``文化仲裁者''著称：他一方面以拉丁基督教国王身份参与十字军时代的地中海政治，另一方面在宫廷中延用阿拉伯语行政传统与伊斯兰世界工匠、学者网络。建筑因而成为最具说服力的合法性舞台——石材、马赛克、木作天花与纺织铭文共同回答``谁统治西西里''这一问题。

\begin{figure}[htbp]
  \centering
  \reportfig{0.85\textwidth}{image5.jpeg}
  \caption{帕拉蒂纳礼拜堂（\emph{Cappella Palatina}）内景：拜占庭金色马赛克与伊斯兰式穆喀纳斯（\emph{muqarnas}）木质天花并置，呈现三大文明在同一宗教空间中的层叠。}
  \label{fig:cappella}
\end{figure}

诺曼宫（图\ref{fig:palace}）在原有阿拉伯宫廷格局上改建扩展：外立面以厚重体量、防御性塔楼与连拱券展现诺曼王权，却保留伊斯兰建筑中常见的盲拱与几何韵律\cite{hillenbrand1999}。位于宫内的帕拉蒂纳礼拜堂（图\ref{fig:cappella}）则是合成文化最集中的空间：外部结构与入口保留诺曼式特征，内部东端圣所覆盖拜占庭风格金色马赛克，主穹顶为``全能基督''，周围环列天使、先知与圣徒；中殿上方却悬置整片绘有宴饮、乐舞、狩猎、动物纹样与祈福性阿拉伯铭文的穆喀纳斯（\emph{muqarnas}）木质天花\cite{fein_cappella,hillenbrand1999}。同一宗教建筑内，基督教圣像叙事与伊斯兰式立体装饰并置，在12世纪地中海极为罕见。诺里奇在《地中海史》中评价，罗杰二世最惊人的政治成就之一，是把拉丁、希腊与阿拉伯三种文明视觉传统集于同一礼拜空间，而与此同时，这三种文明在欧洲大陆许多地区仍彼此征战\cite{norwich2011}。

\begin{figure}[htbp]
  \centering
  \reportfig{0.75\textwidth}{image3.jpeg}
  \caption{帕拉蒂纳礼拜堂穆喀纳斯天花细部：八角星形格室与细密彩绘，体现伊斯兰几何装饰与诺曼宫廷审美需求。}
  \label{fig:muqarnas}
\end{figure}

图\ref{fig:muqarnas}所示天花细部表明，穆喀纳斯不仅是装饰，更是几何秩序与宫廷审美的技术结晶：八角星形格室、细密彩绘与层叠凹凸结构，将伊斯兰建筑中的``无限延伸''感引入基督教礼拜空间。与此相应，罗杰二世加冕斗篷以库法体阿拉伯铭文、丝绸、金线刺绣、珍珠与宝石装饰，在巴勒莫王室工坊制作（伊斯兰历528年/公元1133--1134年），与伊斯兰 \emph{tiraz} 宫廷纺织传统相呼应\cite{schmitz_mantle,dolezalek2017}。Dolezalek 指出，基督教国王身着带阿拉伯文字的礼服，并非``伊斯兰化''，而是利用跨文化符号展示财富、技艺与地中海霸权。诺曼国王借阿拉伯文字与纺织工艺宣示王权——阿拉伯影响已深入建筑空间、宗教图像、天花与服饰等多重载体。除帕拉蒂纳礼拜堂外，巴勒莫地区的齐萨宫（\emph{Zisa}）等亦保留阿拉伯式庭院、喷泉与灰泥装饰，说明伊斯兰空间美学在诺曼时期被选择性保留，服务于国王的奢华政治与外交接待。值得注意的是，这些建筑并非简单``拼贴''，而是按照诺曼宫廷礼仪重新编排：基督教礼拜、伊斯兰式休闲空间与拜占庭式神圣图像各司其职，共同支撑罗杰王朝的统治想象。艺术史因而应被读作政治史：每一类装饰都在回答谁有资格定义``美''与``神圣''。

\section{日常生活与市场网络}

\begin{figure}[htbp]
  \centering
  \reportfig{0.72\textwidth}{image6.jpeg}
  \caption{巴勒莫传统市集街巷（如 Ballarò 一带）：街头贸易与饮食习俗延续数百年跨文化交换的成果。}
  \label{fig:market}
\end{figure}

若将阿拉伯影响仅理解为宫殿艺术，便忽视其进入饮食、农业与市集的方式。9世纪后，北非与东方地中海的农业知识随征服与贸易进入西西里：灌溉渠系、园艺管理与新作物（如柑橘、茄子、菠菜、棉花、硬粒小麦等）逐步改变岛上饮食结构\cite{metcalfe2009,finley1986}。这些变化并非停留在文献记载，而沉淀为可感知的日常实践——巴勒莫传统市集（如 Ballarò 一带，见图\ref{fig:market}）中，水果、香料、鱼类与街头小吃的组合，仍可见地中海跨文化交换的长期后果。市集是``微观史''的现场：商贩叫卖、顾客议价与狭窄街巷中的气味与声音，把数百年前的作物传播与港口贸易转化为当代生活的背景。

诺曼统治时期，巴勒莫仍是多族群、多语言、多宗教并存的城市；希腊人、阿拉伯人、犹太人、拉丁人与诺曼人共同参与贸易、手工与文书工作\cite{johns2002,metcalfe2009}。Johns 对皇家迪旺的研究表明，诺曼国家机器在很大程度上依赖阿拉伯语行政传统，使穆斯林书吏与法官在征服后仍长期担任关键职务。这种共存并不等同于现代意义上的平等，冲突与歧视亦史有所载；但它说明西西里文明并非由单一民族``创造''，而是在接触、协商与再创造中生成。混合文明的韧性最终体现在厨房与街巷：宫殿可焚毁，王朝可更替，主食结构、贸易网络与方言借词却可延续数百年。对课程而言，这一层面回应了``阿拉伯文化如何进入并留在西西里''：答案不仅在宫殿，更在田地、餐桌与市集。近代以来，西西里又经历西班牙哈布斯堡、波旁王朝与意大利统一等政权更替，但地方饮食与地名中的阿拉伯借词仍可作为历史层累的证据。研究西西里，因而也是研究阿拉伯世界影响如何穿越政治边界、渗入民间社会的长时段过程。

\section{知识翻译与易德里西地图}

诺曼西西里是阿拉伯、希腊与拉丁世界的十字路口。11--13世纪，欧洲对古典知识的需求上升，翻译运动成为连接伊斯兰学术遗产与拉丁基督教学术的关键机制。西班牙托莱多在12世纪以阿拉伯语—拉丁语翻译著称，但西西里路径有其独特之处：宫廷能\textbf{同时}接触希腊语原文与阿拉伯语中继文本，形成``双路径''知识输入\cite{molinini2009,metcalfe2009}。这意味着诺曼统治者不必完全依赖伊比利亚或中东的中间环节，便可在巴勒莫组织多语种学术生产。

\begin{figure}[htbp]
  \centering
  \reportfig{0.88\textwidth}{image7.jpeg}
  \caption{易德里西（\emph{al-Idrīsī}）《罗杰之书》（1154）世界地图局部（\emph{Tabula Rogeriana}）：采用阿拉伯传统上南下北朝向，分气候带呈现已知世界。}
  \label{fig:idrisi}
\end{figure}

约1138年，阿拉伯地理学家易德里西（\emph{al-Idrīsī}）应罗杰二世之邀来到巴勒莫，1154年完成《罗杰之书》（\emph{Kitāb Rujār} / \emph{Tabula Rogeriana}）：将已知世界划分为七个气候带，每带再分十区，共70幅分区图及一幅圆形世界地图（图\ref{fig:idrisi}），并采用阿拉伯制图传统上南下北的朝向\cite{charta_rogeriana,metcalfe2009,houben2002}。罗杰二世还出资用约466两白银打造银制天球仪，以配合地理著作的展示与教学。一位拉丁基督教国王邀请穆斯林学者完成世界级地理工程，其象征意义远超个人好学：它表明诺曼政权把伊斯兰学术视为国家声望与文化资本的一部分。易德里西并非``臣服''于欧洲君主，而是在宫廷庇护下与统治者共同生产知识产品。

\begin{figure}[htbp]
  \centering
  \reportfig{0.88\textwidth}{image8.jpeg}
  \caption{大英图书馆藏哈雷三语诗篇（\emph{Harley MS 5786}）书页：希腊语、拉丁语、阿拉伯语三栏对照，红墨水连线标示句级对应，是宫廷多语协作的物质证据。}
  \label{fig:psalter}
\end{figure}

罗杰二世与威廉一世时期，巴勒莫宫廷将托勒密、柏拉图、亚里士多德等著作系统译入拉丁语；翻译活动依赖阿拉伯语、希腊语与拉丁语书吏的协作，形成所谓``西西里翻译学派''\cite{molinini2009}。大英图书馆藏哈雷三语诗篇（\emph{Harley MS 5786}，见图\ref{fig:psalter}）是这一协作的物质证据：一页羊皮纸上希腊语、拉丁语与阿拉伯语三栏并列，红墨水横线标示句级对应；希腊栏与拉丁栏之间存在同步勘误痕迹，说明书吏在同一空间内即时互校\cite{ohogan2022}。1224年，罗杰二世之孙、神圣罗马帝国皇帝腓特烈二世创立那不勒斯大学，将阿拉伯文手稿与译本纳入高等教育体系\cite{metcalfe2009}，使西西里经验成为拉丁欧洲知识革命的重要前奏之一。与课程讨论的奥斯曼``笔墨治国''相呼应：西西里案例显示，跨文明翻译同样是帝国治理与文化竞争的核心技术。托莱多学派擅长把阿拉伯哲学与科学文本转译为拉丁语，却常失去希腊语原典的独立通道；西西里则因长期存在希腊语社群与拜占庭文化遗产，得以在翻译链上保留``希腊语—阿拉伯语—拉丁语''的三重校验。这一结构优势，使诺曼西西里在中世纪欧洲知识史中的地位不应被低估。

\section{结论}

综观上述案例，可归纳四点结构性认识。第一，文化影响力的持久性不必依赖军事统治的连续：阿拉伯军队在11世纪末退出西西里政治舞台，但灌溉、作物、迪旺行政、建筑装饰与宫廷礼仪在诺曼时期反而更系统化\cite{johns2002,metcalfe2009}，反驳``征服即文化断裂''的简化史观。第二，公共媒介体现不同的``可见性''规训：威尼斯杜卡特以圣像与拉丁铭文铸入神学—政治宣言\cite{stahl2000,grierson1990}，伊斯兰金币以书法与名号回避人像\cite{grierson1960,hillenbrand1999}，诺曼西西里则通过 \emph{tarì} 记忆、宫廷纺织与建筑合成居于两者之间。第三，宫廷制度化宽容是跨文化知识生产的关键：易德里西地图与哈雷三语诗篇表明，知识传播的最优路径常是资助者、译者与学者的合作，而非征服者对文献的单向搬运\cite{molinini2009,ohogan2022,houben2002}。第四，混合文明的韧性最终体现在日常生活：市集、饮食、语言混用与行政惯例构成深层基底，使西西里在政权更替后仍保持可识别的``阿拉伯—地中海''底色。

回到课程《阿拉伯世界的1500年》的核心关切：阿拉伯文明并非封闭于东岸与北非，而是与希腊、罗马、拜占庭、拉丁及后来的突厥—波斯传统持续对话的流动体。西西里案例提示，地中海历史的主轴往往不是``纯粹''的征服或皈依，而是如帕拉蒂纳礼拜堂穆喀纳斯天花与拜占庭马赛克并置那般——不同传统在同一空间、同一枚钱币、同一页羊皮纸上相互校正、共存与新生。希提在讨论帝国兴衰时强调，长治久安靠的不仅仅是刀剑，更是笔墨、称号与象征资源的制度化运用\cite{hitip1958}；诺曼西西里以迪旺档案、易德里西地图与罗杰二世斗篷给出了在地验证。对于今日读者而言，这一历史遗产的意义在于：文化边界从来不是铁板一块，制度与符号的创造性重组，往往是多元社会得以延续的关键。将西西里置于更长的阿拉伯—地中海史中观察，还可看到伊斯兰文化在伊比利亚、南意大利与马格里布之间形成的环形交流带；西西里是该带上向北进入拉丁欧洲的重要节点。研究这一节点，有助于摆脱把阿拉伯文明仅视为``东方他者''的视角，转而理解其作为中世纪欧洲形成过程中内在要素的角色。本文材料主要依据课程所论及之希提《阿拉伯通史》相关章节，并参酌近年中世纪地中海与西西里研究的新成果\cite{hitip1958,alshaar2026}。

\clearpage
\renewcommand{\refname}{参考文献}
\begin{thebibliography}{99}
\setlength{\itemsep}{0pt}

\bibitem{johns2002}
J.~Johns, \emph{Arabic Administration in Norman Sicily: The Royal Diwan}. Cambridge, U.K.: Cambridge Univ. Press, 2002.

\bibitem{metcalfe2009}
A.~Metcalfe, \emph{The Muslims of Medieval Italy}. Edinburgh, U.K.: Edinburgh Univ. Press, 2009.

\bibitem{grierson1990}
P.~Grierson, \emph{Money and Its Use in Medieval Europe, 800--1500}. Cambridge, U.K.: Cambridge Univ. Press, 1990.

\bibitem{stahl2000}
A.~M.~Stahl, \emph{Zecca: The Mint of Venice in the Middle Ages}. Baltimore, MD, USA: Johns Hopkins Univ. Press, 2000.

\bibitem{travaini2003}
L.~Travaini, ``Saints and sinners: Coins and religions in medieval Europe,'' \emph{The Numismatic Chronicle}, vol.~163, pp.~177--190, 2003.

\bibitem{grierson1960}
P.~Grierson, \emph{The Coinage of Islam in the Seventh Century}. Oxford, U.K.: Oxford Univ. Press, 1960.

\bibitem{heidemann2010}
S.~Heidemann, ``The coins of the early Islamic period,'' in \emph{The New Cambridge History of Islam}, vol.~1. Cambridge, U.K.: Cambridge Univ. Press, 2010.

\bibitem{hillenbrand1999}
R.~Hillenbrand, \emph{Islamic Art and Architecture}. London, U.K.: Thames \& Hudson, 1999.

\bibitem{metcalf1998}
D.~M.~Metcalf, ``Arabic coins in Norman Sicily,'' in \emph{The Society of Norman Italy}. Leiden, Netherlands: Brill, 1998.

\bibitem{loud2012}
G.~A.~Loud, \emph{Roger II and the Creation of the Kingdom of Sicily}. Manchester, U.K.: Manchester Univ. Press, 2012.

\bibitem{sicilygeo2026}
查字典地理网, ``西西里岛,'' 地理百科. [Online]. Available: \url{https://dili.chazidian.com/baike-2139/}. Accessed: May 26, 2026.

\bibitem{finley1986}
M.~I.~Finley, \emph{A History of Sicily}. London, U.K.: Chatto \& Windus, 1986.

\bibitem{alshaar2026}
AramcoWorld, ``Book review: Muslim Sicily by Nuha Alshaar.'' [Online]. Available: \url{https://www.aramcoworld.com/en/resources/reviews/2026/ma26/muslim-sicily-encounters-and-legacy}. Accessed: May 26, 2026.

\bibitem{fein_cappella}
A.~Fein, ``The Cappella Palatina,'' Smarthistory. [Online]. Available: \url{https://smarthistory.org/cappella-palatina/}. Accessed: May 26, 2026.

\bibitem{norwich2011}
约翰·朱利叶斯·诺里奇, \emph{地中海史：上}, 殷亚平等, 译. 上海, 中国: 东方出版中心, 2011.

\bibitem{schmitz_mantle}
K.~Schmitz von Ledebur, ``Coronation mantle,'' Discover Islamic Art. [Online]. Available: \url{https://islamicart.museumwnf.org/database_item.php?id=object;EPM;at;Mus22;44;en}. Accessed: May 26, 2026.

\bibitem{dolezalek2017}
I.~Dolezalek, \emph{Arabic Script on Christian Kings: Textile Inscriptions on Royal Garments from Norman Sicily}. Berlin, Germany: De Gruyter, 2017.

\bibitem{molinini2009}
D.~Molinini, ``The first Sicilian school of translators,'' \emph{Nova Tellus}, vol.~27, no.~1, pp.~191--205, 2009.

\bibitem{charta_rogeriana}
Univ. of Macau Library, ``Charta Rogeriana.'' [Online]. Available: \url{http://lunamap.must.edu.mo/luna/servlet/detail/MUST~2~2~581~720}. Accessed: May 26, 2026.

\bibitem{houben2002}
H.~Houben, \emph{Roger II of Sicily: A Ruler between East and West}. Cambridge, U.K.: Cambridge Univ. Press, 2002.

\bibitem{ohogan2022}
C.~O{\textquotesingle}Hogan, ``The Harley trilingual psalter,'' in \emph{Medieval Multilingual Manuscripts}. Berlin, Germany: De Gruyter, 2022, pp.~165--182.

\bibitem{hitip1958}
菲利普·希提, \emph{阿拉伯通史}. 北京, 中国: 商务印书馆, 1958.

\end{thebibliography}

\end{document}
